%% =====================================================================
%%  EDIA STATEMENT  (Equity, Diversity, Inclusion and Accessibility)
%% =====================================================================
%%  Institutions use different names for this document. Change the
%%  title in braces below to the one in the job posting, e.g.
%%    {Statement on Equity, Diversity and Inclusion}
%%    {Diversity Statement}   {Inclusive Excellence Statement}
%%    {Statement of Contributions to Community and Belonging}
%%  `short' is the name used in the footer and on the cover page.
%%  `page limit' prints a warning if the statement runs over.
%% =====================================================================
\begin{applicationdocument}[
  short      = {EDIA Statement},
  page limit = 2,
]{Equity, Diversity, Inclusion and Accessibility Statement}

\tip[Structure that works]{Most guides recommend three parts:
(1)~your understanding of equity, diversity, inclusion and accessibility
and why they matter in your field; (2)~concrete past actions and their
outcomes; (3)~specific plans for \emph{this} institution, ideally tied
to its published EDI plan. Aim for one to two pages written in the first
person. You do not have to disclose any part of your identity. Expectations
vary by country and institution: some institutions have made these
statements optional or no longer request them. Write one only when the
posting asks for it, and use the name the posting uses.}

\section{Understanding and Commitment}
Environmental data science shapes decisions about whose water is
monitored, whose knowledge counts as data, and who is trained to do the
work. Each of these is an equity question. Many students, including
students who must work during term, students with disabilities and
first-generation students, never reach the research opportunities that
lead to careers in my field. The barriers are concrete and fixable:
unpaid research positions, inaccessible field and software environments,
and an unspoken curriculum. Indigenous communities have often had their
environmental data collected, used and published without their consent or
benefit. I see inclusive teaching, equitable access to research and
respectful data partnerships as part of doing this science well, not as
work separate from it.

\section{Actions and Outcomes}

\tip[Evidence, not adjectives]{For each example, say what \emph{you}
did, for whom, and what changed, with numbers where you have them. Avoid
``diversity by proxy'', that is, claiming credit for the work of programmes
you only belonged to. It is fine to say what you learned from something
that did not work.}

\textbf{Accessible and transparent teaching.} When I redesigned ENVS~220, I
applied universal design for learning throughout. I wrote captioned,
screen-reader-tested notebooks that run in a browser, so students need
neither a powerful laptop nor any installation. I published explicit
marking criteria for every task, and I introduced flexible deadline tokens.
The share of students who withdrew or failed fell from 18\% to 9\%. The
grade gap between first-generation and continuing-generation students
narrowed from 9 to 3 percentage points.

\textbf{Paid pathways into research.} In 2021 I co-founded \emph{Pathways in
Environmental Data Science}, a paid summer programme for undergraduates
from groups under-represented in the environmental sciences. Over three
summers it hosted 18 students. Eleven of them have since entered graduate
school or environmental careers. Securing stable funding was the
hardest part. I learned to write mentee salaries into grants from the
start, a practice I will continue as a faculty member.

\textbf{Respectful data partnerships.} At Northbridge, I worked with a
First Nation's water monitoring programme to integrate its lake data into
BloomCast. Guided by the community's priorities and the First Nations
principles of OCAP\textsuperscript{\textregistered} (ownership, control,
access and possession), we co-wrote a data agreement. Under it the
community keeps ownership of its data and approves every publication. The
community's monitors now receive forecasts for their lakes first.

\textbf{Service.} As the postdoctoral representative on my faculty's EDI
committee, I helped draft a field-work accessibility protocol, adopted in
2025, that sets out accommodation planning and safety procedures for
remote sites.

\section{Plans at \Institution}

\Institution's EDI action plan commits to removing barriers in
undergraduate research and to strengthening partnerships with Indigenous
communities. My plans support both:
\begin{itemize}
  \item \textbf{Teaching:} bring studio-based, universally designed teaching
        to the \Department's quantitative courses, and share the materials
        openly with colleagues.
  \item \textbf{Lab culture:} publish a lab charter on expectations,
        authorship and mentoring; budget paid undergraduate positions in
        every grant; recruit through the \Faculty's first-generation and
        Indigenous student networks.
  \item \textbf{Research:} pursue partnered research only on terms set
        with communities, and budget for community co-researchers.
  \item \textbf{Service:} contribute to the \Faculty's mentoring programmes
        for first-generation students, and keep assessing my own practice
        against student outcomes.
\end{itemize}

\end{applicationdocument}
